\subsection{Optical Contact Bonding}
Optical contact bonding is a glueless process whereby two closely conformal 
surfaces are joined, being held purely by intermolecular forces. 
\\\\
Generally, intermolecular forces are not strong enough to create
a robust bond between two substrates. However, if the surfaces to be bound
are conformal with nanometric precission, a sufficient surface area is in close 
enough contact ---also known as \textit{intimate contact}--- for the intermolecular
interactions to have an observable macroscopic effect.
\\\\
In practice, this means 
that, in order to bond two glass surfaces together (e.g. two disks to create a wafer),
they must meet very strict roughness and waviness requirements. 
Nowadays, many commercially available glass disks meet this requirements,
\textbf{ESPECIFICAR I POSAR CITA}. In literature, bonding energies up to
\textbf{VALOR} have been achieved. To obtain such strong bonds,
special treatments are often applied to the substrates, in order to ensure
surfaces are very clean and ready to form intermolecular
bonds between them. To do so, many chemical, plasma, pressure and thermal
treatments have proven to be effective.
\\\\
Since Optical Contact Bonding does not require any adhesive, and low-temperature
methods have been developed during the last decades, it is a promising bonding
procedure for glass-based atomic sensors.